\documentclass[11pt]{article}
\usepackage[T1]{fontenc}
\usepackage{lmodern}
\usepackage{amsmath,amssymb,amsthm,mathtools}
\usepackage[margin=1in]{geometry}
\usepackage{microtype}
\usepackage{booktabs}
\usepackage{enumitem}
\usepackage[hidelinks]{hyperref}

\newcommand{\A}{\mathbb A}
\newcommand{\C}{\mathbb C}
\newcommand{\Gm}{\mathbb G_m}
\newcommand{\Pj}{\mathbb P}
\newcommand{\Lef}{\mathbb L}
\newcommand{\Res}{\operatorname{Res}}
\newcommand{\Sym}{\operatorname{Sym}}
\newcommand{\Pic}{\operatorname{Pic}}
\newcommand{\Cl}{\operatorname{Cl}}
\newcommand{\Spec}{\operatorname{Spec}}

\newtheorem{theorem}{Theorem}[section]
\newtheorem{proposition}[theorem]{Proposition}
\newtheorem{lemma}[theorem]{Lemma}
\newtheorem{corollary}[theorem]{Corollary}
\theoremstyle{remark}
\newtheorem{remark}[theorem]{Remark}

\title{Hyperplane Slices of the Marked-Cubic Map\\
\large A classification and geometric explanation of the three-dimensional Keller counterexample}
\author{Prepared by GPT-5.6 Pro (OpenAI)\\
with research direction from Ian Pitchford}
\date{27 July 2026}

\begin{document}
\maketitle

\begin{abstract}
Let $W$ be a two-dimensional complex vector space, put
$V_d=\Sym^d(W^*)$, and consider multiplication of binary forms
\[
m:V_1\times V_2\longrightarrow V_3,\qquad (L,Q)\longmapsto LQ.
\]
For $0\ne\ell\in V_3^*$ define
\[
H_\ell=\{C\in V_3:\ell(C)=1\},\qquad
X_\ell=\{(L,Q):\Res(L,Q)=1,\ \ell(LQ)=1\}.
\]
We prove that $f_\ell:X_\ell\to H_\ell$ is \`etale for every $\ell$ and
classify exactly when $X_\ell$ is affine three-space. If
$\Delta\subset\Pj(V_3)$ is the twisted cubic of cubes of linear forms, then
\[
X_\ell\cong\A^3
\quad\Longleftrightarrow\quad
\Pj(\ker\ell)\text{ is tangent but not osculating to }\Delta.
\]
The tangent case is an iterated Zariski-locally-trivial $\A^1$-bundle. For
contact types $1+1+1$, $2+1$, and $3$, respectively, we compute
\[
[X_\ell]=\Lef^3-\Lef,\qquad \Lef^3,\qquad \Lef^3-\Lef^2.
\]
The osculating slice is $\Gm\times\A^2$. We also identify the fibers with the
simple roots of the target cubic, describe the discriminant as the
nonproperness locus, and reconcile the geometry with the announced explicit
polynomial map. An exact symbolic verifier accompanies the manuscript.
\end{abstract}

\section{Status and scope}

A concrete polynomial self-map of $\C^3$ with nonzero constant Jacobian and a
three-point fiber was announced by Levent Alp\"oge on 19 July 2026 and
credited to Fable. In the following days, Andy Jiang publicized the
symmetric-product marked-root construction; Aaron Lou gave a reproducible
factorization--resultant derivation; and Will Sawin, David Speyer, Daniel
Litt, Jake Levinson, Terence Tao, Bartosz Naskr\k{e}cki, and others supplied
complementary geometric and computational explanations.

This paper does not claim priority for the counterexample or those ingredients.
It assembles them into a complete classification of the affine hyperplanes
$\ell(C)=1$. We work over $\C$. Characteristics two and three require separate
treatment and are not claimed here.

\section{The normalized factorization family}

Let $V_i=\Sym^i(W^*)$. Multiplication
\[
m:V_1\times V_2\to V_3,\qquad (L,Q)\mapsto LQ
\]
has the relative scaling symmetry
\[
(L,Q)\mapsto(\lambda L,\lambda^{-1}Q).
\]
The resultant $\rho(L,Q)=\Res(L,Q)$ is bihomogeneous of degree $(2,1)$:
\[
\rho(\alpha L,\beta Q)=\alpha^2\beta\rho(L,Q).
\]
Thus relative scaling changes $\rho$ by one power of $\lambda$, and $\rho=1$
fixes that scaling uniquely wherever $L$ and $Q$ are coprime.

For $0\ne\ell\in V_3^*$ put
\[
H_\ell=\{C:\ell(C)=1\}\cong\A^3,
\qquad
X_\ell=\{(L,Q):\rho(L,Q)=1,\ \ell(LQ)=1\},
\]
and let $f_\ell(L,Q)=LQ$.

\section{Universal \`etaleness}

\begin{proposition}\label{prop:universal-etale}
The map
\[
\Phi:V_1\times V_2\longrightarrow V_3\times\A^1,
\qquad \Phi(L,Q)=(LQ,\rho(L,Q))
\]
is \`etale on $\rho\ne0$.
\end{proposition}

\begin{proof}
At a coprime pair, a tangent vector in the kernel of multiplication satisfies
\[
\dot LQ+L\dot Q=0.
\]
Coprimality gives $L\mid\dot L$, hence $\dot L=tL$ and $\dot Q=-tQ$.
Therefore
\[
\ker(dm)_{(L,Q)}=\C(L,-Q).
\]
Bihomogeneity gives
\[
d\rho_{(L,Q)}(L,-Q)=(2-1)\rho(L,Q)=\rho(L,Q)\ne0.
\]
The resultant detects the unique tangent direction forgotten by
multiplication. Since both spaces have dimension five, $d\Phi$ is an
isomorphism.
\end{proof}

\begin{corollary}\label{cor:slice-etale}
For every nonzero $\ell$, the map $f_\ell:X_\ell\to H_\ell$ is \`etale.
\end{corollary}

\begin{proof}
It is the base change of Proposition~\ref{prop:universal-etale} along
$H_\ell\times\{1\}\hookrightarrow V_3\times\A^1$.
\end{proof}

\section{Projective and marked-root models}

Let
\[
\overline X=\Pj(V_1)\times\Pj(V_2)\cong\Pj^1\times\Pj^2.
\]
Let $\mathcal R$ be the resultant divisor and
$\mathcal S_\ell=\{([L],[Q]):\ell(LQ)=0\}$. Their divisor classes are
\[
[\mathcal R]=(2,1),\qquad[\mathcal S_\ell]=(1,1).
\]

\begin{lemma}\label{lem:projective-complement}
There is a canonical isomorphism
\[
X_\ell\cong\overline X\setminus(\mathcal R\cup\mathcal S_\ell).
\]
\end{lemma}

\begin{proof}
For a projective pair outside the boundary, choose representatives and put
$r=\rho(L,Q)$ and $h=\ell(LQ)$. The unique normalization is
\[
L'=\frac{h}{r}L,\qquad Q'=\frac{r}{h^2}Q.
\]
Both formulas are regular where $rh\ne0$, and they give
$\rho(L',Q')=\ell(L'Q')=1$.
\end{proof}

A linear form determines its zero $p\in\Pj^1$.

\begin{lemma}\label{lem:marked-root}
There is a natural isomorphism
\[
X_\ell\cong
\{(C,p):C\in H_\ell,\ C(p)=0,\ p\text{ is a simple root of }C\}.
\]
Under this identification, $f_\ell$ forgets the marked root.
\end{lemma}

\begin{proof}
A normalized pair gives $(LQ,Z(L))$. Conversely, a simple root gives a
linear--quadratic factorization, unique up to relative scaling. Since relative
scaling changes the resultant with weight one, there is a unique
resultant-one representative.
\end{proof}

Hence the fiber sizes are $3$, $1$, and $0$ for cubics with three distinct
roots, a double plus a simple root, and a triple root, respectively.

\section{The three hyperplane orbits}

Let
\[
\Delta=\{[M^3]:[M]\in\Pj(V_1)\}\subset\Pj(V_3)
\]
be the twisted cubic. The projective hyperplane
$H_\ell^\infty=\Pj(\ker\ell)$ cuts a degree-three divisor on $\Delta\cong\Pj^1$.
Up to $PGL_2$, the multiplicity types are
\[
1+1+1,\qquad 2+1,\qquad 3.
\]
They are, respectively, transverse, tangent nonosculating, and osculating.

If $\ell'=\lambda\ell$, then
\[
(L,Q)\mapsto(\lambda L,\lambda^{-2}Q)
\]
identifies $X_\ell$ and $X_{\ell'}$. The $SL_2$-equivariance of multiplication
and invariance of the resultant show that the isomorphism class depends only
on the contact type.

\section{The tangent slice as an iterated affine-line bundle}

Assume $H_\ell^\infty\cap\Delta=2\infty+0$. Choose an affine coordinate on
$\Pj^1\setminus\{\infty\}$ so the finite hyperplane condition is
$p+q+r=0$. In $\Sym^2(\Pj^1)\cong\Pj^2$, define
\[
R_p=\{\{q,r\}:p\in\{q,r\}\},
\qquad
S_p=\{\{q,r\}:p+q+r=0\}.
\]
The lines meet at
\[
s(p)=\{p,-2p\}
\]
for finite $p$, and $s(\infty)=\{\infty,\infty\}$. The image $C_s$ of
$s:\Pj^1\to\Pj^2$ is a smooth conic.

Define
\[
X_\ell^*=\{(p,\Lambda):s(p)\in\Lambda,\ \Lambda\ne R_p,S_p\}.
\]

\begin{lemma}\label{lem:first-a1}
The map
\[
\pi:X_\ell\to X_\ell^*,\qquad (p,D)\mapsto(p,\overline{s(p)D})
\]
is a Zariski-locally-trivial $\A^1$-bundle.
\end{lemma}

\begin{proof}
For fixed $(p,\Lambda)$, the point $D$ ranges over
$\Lambda\setminus\{s(p)\}\cong\A^1$. Globally this is the complement of a
section in a $\Pj^1$-bundle.
\end{proof}

A line through $s(p)$ has residual intersection $s(q)$ with $C_s$. The line
$S_p$ joins $s(p)$ to $s(\infty)$, while $R_p$ meets $C_s$ again at
$s(-p/2)$. Therefore
\[
X_\ell^*\cong
\{(p,q)\in\Pj^1\times\A^1:p\ne-2q\}.
\]
Projection to $q$ is the complement of a section in a trivial $\Pj^1$-bundle,
so it is an $\A^1$-bundle over $\A^1$.

\begin{lemma}\label{lem:a1-trivial}
Every Zariski-locally-trivial $\A^1$-bundle over $\A^n$ is trivial.
\end{lemma}

\begin{proof}
Its linear part is a line bundle, classified by $\Pic(\A^n)=0$. Once this is
trivial, its translational part is a $\mathbb G_a$-torsor, classified by
$H^1(\A^n,\mathcal O)=0$.
\end{proof}

Thus $X_\ell^*\cong\A^2$, and Lemmas~\ref{lem:first-a1} and
\ref{lem:a1-trivial} give the main geometric result.

\begin{theorem}\label{thm:tangent-a3}
If $H_\ell^\infty$ is tangent but not osculating to $\Delta$, then
$X_\ell\cong\A^3$.
\end{theorem}

\begin{remark}[Vector-bundle form]
Write the linear factor as $[a:b]$ and the quadratic coefficients as
$[c:d:e]$. The two boundary lines are
\[
bc+ad=0,\qquad b^2c-abd+a^2e=0.
\]
Their moving intersection is the conic section
\[
\sigma([a:b])=[a^2:-ab:-2b^2].
\]
Projection from it is encoded by
\[
0\to\mathcal O(-2)
\xrightarrow{(a^2,-ab,-2b^2)^T}\mathcal O^3
\xrightarrow{\left(\begin{smallmatrix}b&a&0\\0&-2b&a\end{smallmatrix}\right)}
\mathcal O(1)^2\to0.
\]
Thus the quotient $\Pj^1$-bundle is trivial. The boundary lines descend to a
constant section and a degree-one graph, another direct route to the two
successive $\A^1$-bundles.
\end{remark}

\section{The other two orbits}

Let $\Lef=[\A^1]$ in $K_0(\operatorname{Var}_{\C})$. Project the marked-root
model to $p\in\Pj^1$. Usually the equations $C(p)=0$ and $\ell(C)=1$ form an
$\A^2$, and the multiple-root condition removes an $\A^1$. Thus the usual
fiber has class $\Lef^2-\Lef$.

Let $K_p$ be the subspace of cubics vanishing to order at least two at $p$.
Then $K_p^\perp$ is the tangent line to the dual twisted cubic at the evaluation
functional $[e_p]$. Therefore:
\begin{itemize}[leftmargin=2em]
\item in type $1+1+1$, every fiber is usual;
\item in type $2+1$, one exceptional fiber is $\A^2$;
\item in type $3$, one exceptional fiber is empty.
\end{itemize}
Stratifying $\Pj^1$ gives
\[
[X_\ell]=
\begin{cases}
(\Lef+1)(\Lef^2-\Lef)=\Lef^3-\Lef,&1+1+1,\\
\Lef(\Lef^2-\Lef)+\Lef^2=\Lef^3,&2+1,\\
\Lef(\Lef^2-\Lef)=\Lef^3-\Lef^2,&3.
\end{cases}
\]
Compactly supported Euler characteristic sends $\Lef$ to $1$, so the first
and third types have Euler characteristic zero and cannot be $\A^3$.

For type $3$, choose the normal form in which $\ell(C)$ is the leading
coefficient. With
\[
L=aT+bS,\qquad Q=cT^2+dTS+eS^2,
\]
the slice has $ac=1$. Thus $a\in\Gm$, $c=a^{-1}$, and the resultant equation
solves for $e$, leaving $(a,b,d)$ free. Hence
\[
X_\ell\cong\Gm\times\A^2.
\]

\begin{theorem}[Classification]\label{thm:classification}
For $0\ne\ell\in V_3^*$,
\[
X_\ell\cong\A^3
\quad\Longleftrightarrow\quad
H_\ell^\infty\text{ is tangent but not osculating to }\Delta.
\]
Equivalently, the successful projective hyperplanes form the tangent
developable of the dual twisted cubic with the twisted cubic removed.
\end{theorem}

\section{Covering and nonproperness geometry}

Lemma~\ref{lem:marked-root} shows that the fiber is the set of simple roots.
Thus the generic degree is three and the image omits exactly the triple-root
locus in $H_\ell$. Over the discriminant complement, $f_\ell$ is a finite
\`etale three-sheeted covering. At a double-root cubic, the simple-root sheet
remains and the two coalescing markings escape to infinity. At a triple-root
cubic, no marking remains. Consequently the discriminant hypersurface is
exactly the nonproperness locus.

If $C(t)$ has simple root $\zeta$, its unique resultant-one factorization is
\[
L(t)=\frac{t-\zeta}{C'(\zeta)},
\qquad
Q(t)=C'(\zeta)\frac{C(t)}{t-\zeta}.
\]
The escape of a sheet as a root becomes multiple is therefore the divergence
of $1/C'(\zeta)$.

\section{Explicit coordinates}

In the tangent normal form,
\[
X=\{(a,b,c,d,e):a^2e-abd+cb^2=1,\ ad+bc=1\}.
\]
A polynomial isomorphism $\A^3_{a,y,z}\to X$ is
\[
\begin{aligned}
b&=1+ay,\\
c&=1-\frac32ay+a^2z,\\
d&=\frac12y-az+\frac32ay^2-a^2yz,\\
e&=-2z+4y^2-4ayz+3ay^3-2a^2y^2z.
\end{aligned}
\]
One polynomial inverse is
\[
y=2bd-ae,
\qquad
z=2d^2+ce+6bd^2+3bce-\frac92e.
\]
After defining $y$, the second formula may equivalently be written
$z=2yd+cy^2-e/2$.

Dropping the fixed coefficient $ad+bc=1$ from the product gives
\[
G(a,y,z)=(ac,ae+bd,be),
\]
with
\[
\begin{aligned}
G_1&=a-\frac32a^2y+a^3z,\\
G_2&=\frac12y-3az+6ay^2-6a^2yz+\frac92a^2y^3-3a^3y^2z,\\
G_3&=-2z+4y^2-6ayz+7ay^3-6a^2y^2z+3a^2y^4-2a^3y^3z.
\end{aligned}
\]
The exact verifier gives $\det DG=-1$.

Define
\[
A(x,y,z)=\left(x,y,-\frac z2\right),
\qquad
B(u_1,u_2,u_3)=(u_3,2u_2,2u_1).
\]
Then the announced map
\[
\begin{aligned}
F_1&=(1+xy)^3z+y^2(1+xy)(4+3xy),\\
F_2&=y+3x(1+xy)^2z+3xy^2(4+3xy),\\
F_3&=2x-3x^2y-x^3z
\end{aligned}
\]
satisfies $F=B\circ G\circ A$. Hence $\det DF=-2$. Moreover,
\[
F(0,0,-1/4)=F(1,-3/2,13/2)=F(-1,3/2,13/2)=(-1/4,0,0).
\]

\section{The master determinant}

For
\[
L=xT+\beta S,\qquad Q=\gamma T^2+\delta TS+\varepsilon S^2,
\]
put
\[
\begin{aligned}
c_3&=x\gamma,&c_2&=x\delta+\beta\gamma,
&c_1&=x\varepsilon+\beta\delta,\\
c_0&=\beta\varepsilon,
&\rho&=x^2\varepsilon-x\beta\delta+\beta^2\gamma.
\end{aligned}
\]
Then
\[
\det\frac{\partial(c_3,c_2,c_1,c_0,\rho)}
{\partial(x,\beta,\gamma,\delta,\varepsilon)}=-\rho^2.
\]
This is the coordinate shadow of Proposition~\ref{prop:universal-etale}. On the
slice $c_2=\rho=1$, eliminating $(\delta,\varepsilon)$ divides by a triangular
constraint block of determinant $x^3$; the affine chart contributes the
compensating factor $x^3$.

\section{Why degree three is exceptional}

For the degree-$n$ marked-root compactification
$\Pj^1\times\Pj^{n-1}\to\Pj^n$, the analogous ramification and hyperplane
pullback classes are
\[
(n-1,1),\qquad(1,1).
\]
If both boundary divisors are irreducible, divisor localization gives
\[
\Cl(U_n)\cong
\mathbb Z^2/\langle(n-1,1),(1,1)\rangle
\cong\mathbb Z/(n-2)\mathbb Z.
\]
Thus the natural complement cannot be $\A^n$ for $n\ge4$. The boundary
classes are unimodular only for $n=3$.

\section{Conclusion}

Three mechanisms are now separated. Multiplication plus the resultant gives
universal local invertibility; forgetting a marked simple root gives universal
generic multiplicity three; and tangent nonosculating contact gives the unique
affine-space source. Tangency places the two boundary lines on an auxiliary
conic section, and projection from that section produces two successive affine
line bundles. The other hyperplane orbits retain computable obstructions.

\section*{Reproducibility}

The accompanying \texttt{verify\_all.py} checks, over exact rational arithmetic,
the forward and inverse coordinate maps, both Jacobian determinants, the
three-point collision, the five-dimensional determinant, the conic syzygy, and
the target cubic discriminant.

\section*{Contribution and AI disclosure}

Ian Pitchford supplied the research direction and the sequence of questions
that isolated the group action, compactification, hyperplane-orbit problem, and
invariant classification target. GPT-5.6 Pro synthesized the proofs, wrote the
manuscript and verifier, and performed the exact symbolic audit. This is a
post-announcement synthesis and classification, not a historical account of
the discovery.

\begin{thebibliography}{99}

\bibitem{Alpoge2026}
L. Alp\"oge,
\emph{Announcement of an explicit counterexample to the Jacobian conjecture},
19 July 2026.
\url{https://x.com/__alpoge__/status/2079028340955197566}

\bibitem{Jiang2026}
A. Jiang,
\emph{Projective symmetric-product description of the counterexample},
20 July 2026.
\url{https://x.com/davikrehalt/status/2079175065695035442}

\bibitem{Levinson2026}
J. Levinson,
\emph{Claude's counterexample to the Jacobian conjecture}, 22 July 2026.
\url{https://levjake.wordpress.com/wp-content/uploads/2026/07/jacobian-claude-counterexample-writeup.pdf}

\bibitem{Litt2026}
D. Litt,
\emph{Affine-line-bundle explanation of the marked-cubic source},
21 July 2026.
\url{https://x.com/littmath/status/2079353531430289734}

\bibitem{Lou2026}
A. Lou,
\emph{Deriving an explicit polynomial counterexample to the Jacobian conjecture},
20 July 2026.
\url{https://aaronlou.com/jacobian_counterexample_derivation.pdf}

\bibitem{Naskrecki2026}
B. Naskr\k{e}cki,
\emph{Exact audit and structural analysis: A three-dimensional Keller counterexample},
21 July 2026.
\url{https://nasqret.github.io/jacobian-counterexample/book/index.html}

\bibitem{SpeyerBlog2026}
D. Speyer,
\emph{The new counterexample to the Jacobian conjecture},
Secret Blogging Seminar, 20--21 July 2026.
\url{https://sbseminar.wordpress.com/2026/07/20/the-new-counterexample-to-the-jacobian-conjecture/}

\bibitem{Tao2026}
T. Tao,
\emph{A digestion of the Jacobian conjecture counterexample},
21 July 2026.
\url{https://terrytao.wordpress.com/2026/07/21/a-digestion-of-the-jacobian-conjecture-counterexample/}

\bibitem{Speyer2024}
D. E. Speyer,
\emph{Richardson varieties, projected Richardson varieties and positroid varieties},
2024, Lemma 3.5.
\url{https://arxiv.org/abs/2303.04831}

\end{thebibliography}

\end{document}
